\documentclass[12pt]{article}
\usepackage[utf8]{inputenc}
\usepackage{graphicx}
\usepackage{hyperref}
\usepackage{float}
\usepackage{geometry}
\usepackage{booktabs}
\geometry{a4paper, margin=1in}

\title{\textbf{SPE Final Project Report} \\ Microservices-Based Drift-Aware MLOps for Cybersecurity Log Anomaly Detection}
\author{Group Members \\ [Your Name] – [Roll No] \\ [Team Member 2] – [Roll No]}
\date{\today}

\begin{document}

\maketitle

\tableofcontents
\newpage

\section{System Overview}
The system is a cloud-native, microservices-driven MLOps platform designed for cybersecurity log anomaly detection. It provides real-time inference on network traffic logs to detect known and zero-day threats. The architecture integrates a hybrid machine learning model (Random Forest + Autoencoder) with automated drift detection and retraining capabilities. It follows robust DevOps practices encompassing containerization, CI/CD, orchestration, and centralized logging to ensure high availability and continuous deployment.

\subsection{Core Functionalities}
\begin{itemize}
    \item Real-time threat detection using a hybrid ML model.
    \item Automated drift detection based on Anomaly Rate and Feature Distribution (KS-Test).
    \item Zero-downtime hot-swapping of models upon retraining.
    \item Scalable microservices orchestrated via Kubernetes with HPA.
    \item Secure storage using HashiCorp Vault.
    \item Centralized logging and monitoring using the ELK Stack.
    \item Automated CI/CD pipelines via Jenkins.
\end{itemize}

\subsection{High-Level Architecture Flow}
The frontend interface routes requests to the Prediction Service, which evaluates incoming network logs against the hybrid model. Predictions are appended to a shared log file. The Drift Detection Service continuously monitors these logs against baseline statistics. Upon detecting drift, an automated retraining pipeline is triggered, and the new model is hot-loaded into the Prediction Service without downtime.

% Suggestion: Add an architecture diagram showing the user, frontend, prediction service, drift service, ML pipeline, and the DevOps stack.
% \begin{figure}[H]
%     \centering
%     \includegraphics[width=\textwidth]{images/high_level_architecture.png}
%     \caption{High-Level System Architecture and Request Flow}
% \end{figure}

\section{Overall System Architecture}
The system utilizes a microservices architecture to decouple the frontend, prediction engine, drift monitoring, and ML pipeline. The infrastructure is managed using a robust DevOps framework.

\subsection{Backend Microservices Architecture}
The backend consists of the following 4 microservices/components:
\begin{enumerate}
    \item \textbf{Prediction Service (FastAPI):} Serves the ML model, exposing REST APIs for single and batch predictions.
    \item \textbf{Drift Detection Service (FastAPI):} Monitors prediction logs, performs statistical tests, and exposes drift status and reports.
    \item \textbf{Frontend Service (Streamlit):} Provides the interactive dashboard for users to upload logs and view threat classifications.
    \item \textbf{ML Pipeline Service (Offline/Batch):} Handles data downloading, preprocessing, training, and baseline generation.
\end{enumerate}

\section{Technology Stack}
\begin{itemize}
    \item \textbf{Frontend:} Streamlit
    \item \textbf{Backend:} Python, FastAPI, Uvicorn
    \item \textbf{Machine Learning:} Scikit-Learn, Pandas, NumPy, SciPy
    \item \textbf{Version Control:} Git, GitHub
    \item \textbf{CI/CD Automation:} Jenkins, GitHub Webhooks
    \item \textbf{Containerization:} Docker, Docker Compose
    \item \textbf{Configuration Management:} Ansible (Roles and Playbooks)
    \item \textbf{Orchestration \& Scaling:} Kubernetes (K8s), Horizontal Pod Autoscaling (HPA)
    \item \textbf{Monitoring \& Logging:} ELK Stack (Elasticsearch, Logstash, Kibana), Filebeat
    \item \textbf{Security:} HashiCorp Vault (Secrets Management)
\end{itemize}

% Suggestion: Add an image of the technology stack layout.
% \begin{figure}[H]
%     \centering
%     \includegraphics[width=0.8\textwidth]{images/tech_stack.png}
%     \caption{Technology Stack Overview}
% \end{figure}

\section{Frontend Architecture and Implementation}
The Frontend Service is built using Streamlit, offering a reactive UI for cybersecurity analysts. 
Users can upload batch CSV files containing network features (NSL-KDD format) or manually enter feature values. The frontend securely communicates with the Prediction Service's REST APIs.

% Suggestion: Add screenshots of the Streamlit dashboard (e.g., File Upload, Prediction Results).
% \begin{figure}[H]
%     \centering
%     \includegraphics[width=0.9\textwidth]{images/frontend_dashboard.png}
%     \caption{Streamlit Frontend Dashboard}
% \end{figure}

\section{Backend Architecture and Implementation}

\subsection{Prediction Service}
The Prediction Service is a highly concurrent FastAPI application. It loads the \texttt{model\_v1.pkl} and \texttt{preprocessor.pkl} into memory.
\textbf{APIs:}
\begin{itemize}
    \item \texttt{POST /predict}: Accepts 41 features, returns a single prediction (Normal, Confirmed Threat, Novel Threat).
    \item \texttt{POST /predict\_batch}: Accepts a CSV file, returns bulk predictions.
    \item \texttt{POST /reload}: Hot-reloads the ML models in-memory without server restart (Zero-downtime).
\end{itemize}

\subsection{Drift Detection Service}
This service calculates statistical deviations between the live production data and the training baseline.
\textbf{APIs:}
\begin{itemize}
    \item \texttt{GET /health}: Service health check.
    \item \texttt{GET /drift/status}: Returns a quick boolean flag if drift is detected.
    \item \texttt{GET /drift/report}: Returns a detailed JSON report containing KS-Statistic and p-values for all 38 numeric features.
    \item \texttt{POST /drift/retrain}: Triggers the background ML retraining process and invokes the prediction service's reload API.
\end{itemize}

\subsection{ML Pipeline (Training \& Retraining)}
A decoupled pipeline that fetches the NSL-KDD dataset, permanently splits it into initial training (70\%) and reserve (30\%), and trains the Hybrid Anomaly Detector. It exports the model artifacts and \texttt{baseline\_stats.json}.

\section{DevOps and Infrastructure}

\subsection{CI/CD Automation with Jenkins}
Jenkins automates the entire software delivery lifecycle. Any incremental push to the GitHub repository triggers a GitHub Webhook, causing Jenkins to:
\begin{enumerate}
    \item Checkout the latest code.
    \item Run automated tests.
    \item Build Docker images for the 4 microservices.
    \item Push the images to Docker Hub.
    \item Trigger continuous deployment to the Kubernetes cluster.
\end{enumerate}

% Suggestion: Add a screenshot of the Jenkins Pipeline success stages.
% \begin{figure}[H]
%     \centering
%     \includegraphics[width=\textwidth]{images/jenkins_pipeline.png}
%     \caption{Jenkins CI/CD Pipeline Stages}
% \end{figure}

\subsection{Containerization and Configuration Management}
All services are containerized using Docker. Multi-container orchestration is defined via \texttt{Docker Compose} for local environments. For broader infrastructure provisioning, \textbf{Ansible Playbooks} are utilized. Modular Ansible Roles are designed to automatically provision servers, install Docker/K8s dependencies, and configure network settings consistently.

\subsection{Kubernetes Orchestration and Security}
The production environment runs on a Kubernetes (K8s) cluster. 
\begin{itemize}
    \item \textbf{Deployments \& Services:} Each microservice runs as a scalable Deployment exposed via ClusterIP.
    \item \textbf{High Availability:} Horizontal Pod Autoscaling (HPA) dynamically scales the Prediction Service based on CPU/Memory utilization during high traffic spikes.
    \item \textbf{Live Patching:} The FastAPI \texttt{/reload} endpoint ensures live patching of the ML model without dropping active HTTP connections.
    \item \textbf{Security (Vault):} Sensitive credentials (such as Docker Hub tokens and API keys) are securely injected into the K8s pods using HashiCorp Vault.
\end{itemize}

% Suggestion: Add an image of Kubernetes nodes/pods running or the HPA in action.
% \begin{figure}[H]
%     \centering
%     \includegraphics[width=\textwidth]{images/k8s_pods.png}
%     \caption{Kubernetes Pods and HPA configuration}
% \end{figure}

\subsection{Monitoring and Logging (ELK Stack)}
To ensure complete observability, centralized logging is achieved via the ELK Stack:
\begin{itemize}
    \item \textbf{Elasticsearch:} Stores and indexes all prediction logs and system access logs.
    \item \textbf{Logstash / Filebeat:} Captures logs from Kubernetes containers and forwards them.
    \item \textbf{Kibana:} Provides analytical dashboards visualizing threat detection rates, anomaly counts, and system performance metrics over time.
\end{itemize}

% Suggestion: Add a screenshot of the Kibana dashboard showing anomaly logs.
% \begin{figure}[H]
%     \centering
%     \includegraphics[width=\textwidth]{images/kibana_dashboard.png}
%     \caption{Kibana Dashboard visualizing Log Anomalies}
% \end{figure}

\section{Conclusion}
This project successfully integrates advanced Machine Learning algorithms with modern DevOps principles. The system not only detects cyber threats accurately using a Hybrid Model but ensures the model's longevity via autonomous Drift Detection and Retraining. Through Jenkins, Kubernetes, Ansible, Vault, and ELK, the platform achieves robust CI/CD, high availability, security, and deep observability, fulfilling the requirements of a production-grade MLOps architecture.

\end{document}
